\documentclass[11pt,a4paper]{article}
\usepackage[a4paper,margin=27mm]{geometry}
\usepackage[T1]{fontenc}
\usepackage{lmodern,microtype}
\usepackage{amsmath,amssymb,amsthm,mathtools}
\usepackage{booktabs,array}
\usepackage{xcolor}
\usepackage[colorlinks=true,citecolor=blue!45!black,linkcolor=blue!45!black,urlcolor=blue!50!black]{hyperref}
\hypersetup{pdftitle={Unforced Navier--Stokes},pdfauthor={Tony Newton},
  pdfsubject={Compact symmetric stresses, exact coefficient margins and formal verification}}
\usepackage{fancyhdr}
\usepackage{enumitem}
\setlength{\parindent}{1.25em}
\setlength{\parskip}{0.20em}
\setcounter{secnumdepth}{2}
\pagestyle{fancy}
\fancyhf{}
\fancyhead[L]{\footnotesize UNFORCED NAVIER--STOKES}
\fancyhead[R]{\footnotesize T. NEWTON}
\fancyfoot[C]{\thepage}
\renewcommand{\headrulewidth}{0.3pt}
\theoremstyle{plain}
\newtheorem{theorem}{Theorem}[section]
\newtheorem{proposition}[theorem]{Proposition}
\newtheorem{lemma}[theorem]{Lemma}
\newtheorem{corollary}[theorem]{Corollary}
\theoremstyle{definition}
\newtheorem{remark}[theorem]{Remark}
\newcommand{\R}{\mathbb{R}}
\newcommand{\T}{\mathbb{T}^{3}}
\newcommand{\e}{\mathrm{e}}
\newcommand{\supp}{\operatorname{supp}}
\newcommand{\diver}{\operatorname{div}}
\newcommand{\cc}{C_c^{\infty}}
\newcommand{\dd}{\,\mathrm{d}}
\newcommand{\norm}[1]{\lVert #1\rVert}
\title{\vspace{-1.2em}\bfseries\LARGE Unforced Navier--Stokes\\[0.45em]
  \large\mdseries Compact symmetric stresses, exact coefficient margins, and formal verification on the three-torus}
\author{Tony Newton\\\small Newton Astro Labs, United Kingdom}
\date{25 October 2026}

\begin{document}
\maketitle
\thispagestyle{plain}
\begin{abstract}
The absence of external forcing makes the three-dimensional Navier--Stokes
regularity problem substantially more restrictive than a construction with
prescribed smooth forcing. We study exact mathematical prerequisites for a
candidate unforced finite-time singularity on the periodic three-torus.
An explicit compact symmetric-divergence inverse is proved for smooth
axisymmetric meridional sources supported in an annulus, with vanishing
total axial force as the necessary and sufficient compatibility condition.
A separate two-diagonal pressure-neutral construction exhibits a sharp
radial sign obstruction. Finite seven-ray stress algebra yields the exact
coefficient margin $15971/122900$ and sharp componentwise amplification
factor $14$. Lean checks the finite cone algebra, radial pressure-gauge
derivative identity, full physical residual expansion, and conditional
retention of terminal speed. These results do not yet construct the
source-matched, scale-sharp, positive-cone-compatible compact stress or
a genuine zero-force singular solution. The global unforced existence
question remains open.
\end{abstract}

\section{The problem and the mathematical objective}

We work at viscosity one on the unit three-torus
$\T=\R^{3}/\mathbb{Z}^{3}$. For a smooth velocity $u(t,x)\in\R^{3}$
and pressure $p(t,x)\in\R$, the unforced incompressible equations are
\begin{equation}\label{eq:NS}
  \partial_tu+(u\cdot\nabla)u-\Delta u+\nabla p=0,
  \qquad \diver u=0.
\end{equation}
The specific target is to find nonzero smooth divergence-free periodic
initial data and a solution smooth for $0\le t<1$, of finite kinetic
energy on a fundamental cell, whose speed becomes unbounded as $t\uparrow1$.
The source-free condition in \eqref{eq:NS} is essential: a verified
construction with a nonzero smooth force does not establish this target.
The Clay problem statement makes this distinction explicit
\cite{Fefferman2000}. The published OpenAI Lean formalization states
Navier--Stokes breakdown with smooth forcing; its unforced finite-time
result concerns Euler rather than Navier--Stokes \cite{OpenAI2026}.

Our approach isolates small results that are meaningful independently of
the unresolved global theorem. First, a source must admit a compact
symmetric-divergence correction. Second, if that correction must remain
in a prescribed positive stress cone, its components need quantitative
control relative to a density envelope. Third, all terms in the actual
nonlinear physical residual must cancel with a single periodic pressure.
Finally, a singular leading trace can survive only if the correction is
provably subleading. Sections 2--6 establish exact parts of this chain.

The construction of compact symmetric anti-divergence operators under
force and torque compatibility is not new; see, for example, the
support-preserving operator and its explicit moment correction in
\cite[\S3.6 and Appendix A]{Zhao2026}. The direct axisymmetric formula
below is offered as an elementary source-specific coordinate realization,
not as a priority claim for compact symmetric anti-divergence. The new
obstacle for the present programme is the additional shrinking-support,
weighted-seven-ray positivity requirement, which is not supplied by
general existence of a compact tensor.

\section{A compact symmetric inverse in an annular region}

Let $0<a<b$ and $c<d$, and write
$\Omega=(a,b)\times(c,d)$ in meridional cylindrical coordinates $(r,z)$.
All functions in this section are independent of the azimuthal angle
$\theta$. The orthonormal cylindrical frame is $(e_r,e_\theta,e_z)$.
For a symmetric tensor with components
\begin{equation}\label{eq:stress}
  S=K\,e_r\otimes e_r+L\,e_\theta\otimes e_\theta
    +M(e_r\otimes e_z+e_z\otimes e_r)
    +C\,e_z\otimes e_z,
\end{equation}
standard cylindrical differentiation gives
\begin{align}
 (\diver S)_r&=\partial_rK+\frac{K-L}{r}+\partial_zM,
       \label{eq:divr}\\
 (\diver S)_z&=\partial_rM+\frac Mr+\partial_zC,
       \label{eq:divz}\\
 (\diver S)_\theta&=0.\notag
\end{align}
The $1/r$ terms are geometric, not lower-order terms that may be
discarded near a shrinking cylindrical support.

\begin{theorem}[Explicit compact meridional inversion]\label{thm:inverse}
Let $f_r,f_z\in\cc(\Omega)$, and extend
$f=f_r e_r+f_z e_z$ by zero to $\R^{3}$ in cylindrical coordinates.
There is a symmetric $S\in\cc(\R^{3};\operatorname{Sym}_3)$,
axisymmetric and supported in the same open annular cylinder, with
$\diver S=f$, if and only if
\begin{equation}\label{eq:vertical-force}
 \int_a^b\int_c^d r f_z(r,z)\dd z\dd r=0.
\end{equation}
The construction uses no nonlocal spatial projection, but it asserts no
positive-semidefinite or weighted-cone bound on $S$.
\end{theorem}

\begin{proof}
For necessity, integrate the axial component of $\diver S=f$
against the cylindrical Jacobian $r$. Compact support and
\eqref{eq:divz} give
\[
 \int_a^b\int_c^d r f_z\dd z\dd r
 =\int_a^b\int_c^d\bigl[\partial_r(rM)+r\partial_zC\bigr]
   \dd z\dd r=0.
\]
For sufficiency, define the axial source marginal and its radial
primitive by
\begin{equation}\label{eq:HJ}
 J(r)=\int_c^d f_z(r,z)\dd z,
 \qquad H(r)=\frac1r\int_a^r sJ(s)\dd s.
\end{equation}
Since $r>0$ on $[a,b]$, $H$ is smooth. The vanishing moment
\eqref{eq:vertical-force} makes $H$ identically zero outside a compact
subinterval of $(a,b)$. Choose any $\eta\in\cc((c,d))$ normalized by
$\int_c^d\eta(z)\dd z=1$, and set
\begin{equation}\label{eq:M}
 M(r,z)=H(r)\eta(z).
\end{equation}
The identity $(rH)'=rJ$ implies
\[
 \int_c^d\left(f_z-\partial_rM-\frac Mr\right)\dd z
   =J-H'-\frac Hr=0.
\]
Consequently the axial primitive
\begin{equation}\label{eq:C}
 C(r,z)=\int_c^z\left[f_z(r,\zeta)-\partial_rM(r,\zeta)
                  -\frac{M(r,\zeta)}r\right]\dd\zeta
\end{equation}
vanishes both below and above the axial support; it is smooth and
compactly supported in $\Omega$. Choose any $K\in\cc(\Omega)$,
including $K=0$, and define
\begin{equation}\label{eq:L}
 L=K+r\bigl(\partial_rK+\partial_zM-f_r\bigr).
\end{equation}
Each component of \eqref{eq:stress} is smooth and compactly supported
away from the axis. Substitution in \eqref{eq:divr} gives
\[
 \partial_rK+\frac{K-L}{r}+\partial_zM
 =\partial_rK-(\partial_rK+\partial_zM-f_r)+\partial_zM=f_r,
\]
and \eqref{eq:C} gives $\partial_rM+M/r+\partial_zC=f_z$.
Thus $\diver S=f$ as claimed.
\end{proof}

\begin{remark}[Scope of the construction]
The theorem concerns a meridional source with $f_\theta=0$ and support
strictly separated from $r=0$. An azimuthal source introduces a
different angular-momentum compatibility condition, and extending
the construction across the axis requires separate regularity
conditions on the cylindrical components. The tensor in the theorem
may be indefinite, and its size can deteriorate when the support
shrinks. None of these omissions is harmless for the weighted H8
attachment problem.
\end{remark}

\section{Pressure-neutral gauges and a sign obstruction}

The mixed and hoop degrees of freedom in Theorem~\ref{thm:inverse} are
not optional decorations. A restricted diagonal ansatz exhibits a
genuine obstruction to maintaining one sign.

\begin{proposition}[Two-diagonal gauge]\label{prop:gauge}
Let $F\in\cc(\Omega)$, put $K=rF_r$ and $C=F+rF_r$, and take
$S=K e_r\otimes e_r+C e_z\otimes e_z$. Then
\begin{equation}\label{eq:gauge}
 \diver S=\nabla C.
\end{equation}
In particular its divergence is a pressure gradient, so it cannot by
itself represent a prescribed source with a nonzero solenoidal part.
\end{proposition}

\begin{proof}
The radial divergence is $K_r+K/r$. Since $K=rF_r$,
$C_r=F_r+K_r=K/r+K_r$. The axial divergence is $C_z$,
which proves \eqref{eq:gauge}.
\end{proof}

\begin{proposition}[Restricted radial sign condition]\label{prop:sign}
If $K$ and $C$ are compactly supported in $r\in(a,b)$ at a fixed $z$
and obey $C_r=K_r+K/r$, then
\begin{equation}\label{eq:momentK}
 \int_a^b K(r,z)/r\dd r=0.
\end{equation}
Consequently, $K\ge0$ forces $K\equiv0$ for this restricted two-diagonal
class.
\end{proposition}

\begin{proof}
Integrate the derivative identity between $a$ and $b$;
the endpoint terms of $K$ and $C$ vanish by compact support.
Since $r>0$, continuity and \eqref{eq:momentK} imply that a
nonnegative $K$ is identically zero.
\end{proof}

This conclusion does not exclude a sign-changing correction around a
positive background tensor or a compact construction with nonzero
hoop and mixed components. Indeed, Theorem~\ref{thm:inverse} explicitly
uses both of those components. The more difficult question is whether
one can choose them while retaining a prescribed strict coefficient
cone with uniform shrinking-support estimates.

\section{The exact seven-ray coefficient margin}

The finite cone appearing in the current H8 source has seven rays.
With aperture $t=1/4$ they are
\begin{equation}\label{eq:rays}
\begin{aligned}
 v_1&=(1,0,0), &v_2&=(1,t,0),&v_3&=(1,-t,0),\\
 v_4&=(1,0,t),&v_5&=(1,0,-t),\\
 v_6&=(1,t,t),&v_7&=(1,t,-t).
\end{aligned}
\end{equation}
Their normalized positive coefficient centre is
$c^{(0)}_i=236/1229$ for $i=3$ and $172/1229$ for every other
index. The raw source covariance and its trace-free rank-one identity
are established by exact finite arithmetic and in
\texttt{UnforcedPatching.SevenRayCenter.lean}.

Write $(A,B,C,D)\in\R^4$ for the four trace-free stress coordinates.
The resulting coefficient perturbations, in ray order, are
\begin{equation}\label{eq:linear-map}
 \delta c=\left(
 -\tfrac32 D,\ 2C-4D,\ -2C-4D,\
 2A-8B+4D,\ -2A+8B+4D,\ 8B,\ -8B\right).
\end{equation}

\begin{proposition}[Sharp componentwise cone estimate]\label{prop:14}
For $m=\max(|A|,|B|,|C|,|D|)$ the map
\eqref{eq:linear-map} satisfies
\begin{equation}\label{eq:14}
 \max_{1\le i\le7}|\delta c_i|\le14m,
\end{equation}
and the number $14$ is optimal. If $\rho\ge0$ and
\begin{equation}\label{eq:eps}
 m\le\frac{15971}{1720600}\rho,
\end{equation}
then for every ray $c_i^{(0)}\rho+\delta c_i\ge\rho/100$.
\end{proposition}

\begin{proof}
The absolute row sums in \eqref{eq:linear-map} are respectively
$3/2$, $6$, $6$, $14$, $14$, $8$, $8$; hence the operator norm
from $\ell^\infty(\R^4)$ to $\ell^\infty(\R^7)$ is at most $14$.
At $(A,B,C,D)=(1,-1,0,1)$, the fourth coordinate equals $14$,
so the constant is sharp. Since
\[
 \frac{172}{1229}-\frac1{100}=\frac{15971}{122900},
\]
inequality \eqref{eq:eps} implies
$|\delta c_i|\le 14m\le(15971/122900)\rho$.
As $c_i^{(0)}\ge172/1229$, the claimed floor follows.
\end{proof}

The seven-ray claim is a finite linear-algebra statement. It provides
the exact smallness target for a potential compact correction but gives
no estimate for the particular tensor delivered by
Theorem~\ref{thm:inverse}. In particular, no implication from
unweighted compact solvability to \eqref{eq:eps} is available.

\section{The physical residual and terminal speed}

Define the viscosity-one residual
\begin{equation}\label{eq:residual}
 \mathcal N(u,p)=u_t+(u\cdot\nabla)u-\Delta u+\nabla p.
\end{equation}
For smooth $(u,p)$ and $(w,q)$, direct expansion yields
\begin{equation}\label{eq:add}
 \mathcal N(u+w,p+q)=\mathcal N(u,p)+\mathcal N(w,q)
   +(u\cdot\nabla)w+(w\cdot\nabla)u.
\end{equation}
Both nonlinear cross terms are essential. Equation \eqref{eq:add},
with explicit differentiability assumptions in the upstream physical
definitions, is checked in Lean by the declaration
\texttt{UnforcedPatching.}\allowbreak
\texttt{navierStokesResidual\_add}.

\begin{proposition}[Conditional terminal retention]\label{prop:ret}
Suppose that $u$ has arbitrarily large speed at points with
$t\uparrow1$. If a correction $w$ satisfies
\begin{equation}\label{eq:subleading}
 \norm{w(t,x)}\le\alpha\norm{u(t,x)}+C_0,
 \qquad 0\le\alpha<1,\quad C_0\ge0,
\end{equation}
uniformly on $0<t<1$ and $x\in\T$, then $u+w$ also has
arbitrarily large speed arbitrarily close to time one.
\end{proposition}

\begin{proof}
The reverse triangle inequality and \eqref{eq:subleading} give
\[
 \norm{u+w}\ge\norm{u}-\norm{w}
 \ge(1-\alpha)\norm{u}-C_0.
\]
Given any target speed $M>0$ and any terminal neighbourhood,
choose a point therein where
$\norm{u}>(M+C_0)/(1-\alpha)$.
Then $\norm{u+w}>M$ at that point.
\end{proof}

This last implication is kernel-checked by the Lean theorem
\begin{center}
\small\texttt{UnforcedPatching.speedUnbounded\_add\_subleading}
\end{center}
It is deliberately conditional: no field $u$ meeting the singularity
premise, and no actual PDE correction $w$ meeting
\eqref{eq:subleading}, has been supplied by these lemmas.

\section{What the formal development proves}

The Lean project is pinned to Lean $4.34.0$-rc2 and a fixed Mathlib
revision. The checked physical interface states smoothness, three
coordinate periods, zero forcing, nonzero initial data, finite
fundamental-cell kinetic energy and endpoint speed divergence. Its
existential proposition is named
\texttt{UnforcedPatching.unforcedTargetStatement}; there is no proof
term inhabiting that proposition.

Actual Lean theorems cover the full nonlinear residual addition
\eqref{eq:add}, divergence and periodicity preservation for sums,
time-only pressure-gauge invariance, a sufficient physical correction
identity, the finite seven-ray data and \eqref{eq:14},
the derivative identity used in Proposition~\ref{prop:gauge}, and
Proposition~\ref{prop:ret}. The printed axioms for these declarations
are the standard logical axioms
\texttt{propext}, \texttt{Classical.choice}, and \texttt{Quot.sound}.
Their scope does not include the cylindrical tensor calculus
\eqref{eq:divr}--\eqref{eq:divz}, the integral construction
\eqref{eq:HJ}--\eqref{eq:L}, compact-support theorems or weighted H8
bounds. Theorem~\ref{thm:inverse} therefore has the complete analytic
proof given above but is \emph{not yet Lean-formalized}.

An independent exact-arithmetic test evaluates a nontrivial annular
polynomial instance of the radial compatibility identity. For the
uncompensated purely radial version, the curl obstruction at
$(r,z)=(5/4,1/2)$ equals $-5103/32768$ exactly; the test explicitly
rejects the pressure-neutral interpretation in this case. The
polynomial is used only for interior differentiation, not as an
infinitely smooth zero-extended cutoff.

The source code, machine-readable G0--G8 ledger, smoke tests and
reproduction commands accompany the public integration branch
\cite{NewtonRepo2026}. A local pinned build has passed ten smoke
modules; at the time of preparing this edition, the corresponding
GitHub-hosted Actions checks had not furnished a successful receipt
for the latest commit. Local compilation does not replace hosted CI,
semantic theorem auditing or independent scrutiny of the paper proof.

\section{Remaining analytic requirements}

The compact meridional inverse of Theorem~\ref{thm:inverse} is a
properly restricted solvability result. It does not resolve the
weighted attachment problem because the desired source may have an
azimuthal component, a noncompact Leray-projected tail, nontrivial
pressure coupling, or incompatible moments. Even when compact
inversion is possible, the output may leave the required seven-ray
positive cone, especially as the support shrinks.

A complete unforced singularity construction must still provide a
single authenticated local profile, source-matched smooth shell
data, a scale-sharp weighted symmetric inverse respecting the
coefficient margin \eqref{eq:eps}, a periodic nonlinear correction
with a proved right-inverse/contraction scheme, the full physical
zero-residual identity, and terminal retention for that actual
correction. Each claim requires a further source-to-Lean semantic
audit and a compiled theorem using the genuine PDE definitions.

Thus the present mathematical result is an exact compact stress
construction and several independently checked necessary or
conditional ingredients. The existence or nonexistence of a
finite-time singular solution to \eqref{eq:NS} with smooth nonzero
periodic initial data is \emph{not decided by this manuscript}.

\begin{thebibliography}{9}
\bibitem{Fefferman2000}
C. Fefferman, \emph{Existence and Smoothness of the Navier--Stokes
Equation}, Clay Mathematics Institute Millennium Prize Problem
statement, \url{https://www.claymath.org/wp-content/uploads/2022/06/navierstokes.pdf}.

\bibitem{OpenAI2026}
OpenAI, \emph{Finite time blowup for Navier--Stokes} (2026), and the
associated formalization repository,
\url{https://github.com/openai/NavierStokesAndEuler}.
The Navier--Stokes result uses smooth forcing; the unforced result
in that repository concerns Euler.

\bibitem{Zhao2026}
X. Zhao, \emph{Compactly Supported Finite-Energy Stationary Solutions of
the Three-Dimensional Navier--Stokes Equations} (2026),
arXiv:2608.17383, \url{https://arxiv.org/abs/2608.17383}.
See the compact symmetric anti-divergence construction in Section 3.6
and Appendix A. Its stationary solutions are distributional; they do
not settle classical smooth-solution regularity.

\bibitem{Leray1934}
J. Leray, \emph{Sur le mouvement d'un liquide visqueux emplissant
l'espace}, Acta Mathematica \textbf{63} (1934), 193--248.

\bibitem{CKN1982}
L. Caffarelli, R. Kohn, and L. Nirenberg, \emph{Partial regularity
of suitable weak solutions of the Navier--Stokes equations},
Communications on Pure and Applied Mathematics \textbf{35} (1982),
771--831.

\bibitem{NewtonRepo2026}
T. Newton, \emph{Unforced Navier--Stokes Lean integration branch},
\url{https://github.com/tonynewton79-web/unforced-navier-stokes-lean/pull/3}
(accessed 25 October 2026).
\end{thebibliography}
\end{document}
